\chapter{Results and Discussion}\label{ch:C}

\section{Dataset Statistics After Augmentation}

The augmentation pipeline expanded the dataset from 529 to \textbf{676 rows}.
Table~\ref{tab:augmented_summary} shows the composition by data source.

\begin{table}[htb]
\centering
\caption{Dataset composition by source after augmentation.}
\label{tab:augmented_summary}
\begin{tabular}{lrl}
\toprule
Data source & Rows & Description \\
\midrule
\texttt{original}  &  13 & All key fields present in the published paper \\
\texttt{estimated} & 516 & STH / H\textsubscript{2} / Bandgap / pH filled via physics model \\
\texttt{synthetic} & 147 & New rows based on literature performance ranges \\
\midrule
Total              & 676 & \\
\bottomrule
\end{tabular}
\end{table}

After removing 33 photocathode rows from the photocurrent target
(Section~\ref{sec:leakage}), the effective training set sizes are:
521 rows for Photocurrent, 556 for STH Efficiency, and 559 for
H\textsubscript{2} Evolution Rate.

The five most represented materials in the dataset are:
BiVO\textsubscript{4} (152 rows), $\alpha$-Fe\textsubscript{2}O\textsubscript{3}
(78 rows), ZnO (40 rows), TiO\textsubscript{2} (36 rows), and WO\textsubscript{3}
(28 rows).
This distribution reflects the prevalence of these materials in published
PEC research over the past two decades.

\section{Model Performance}

\subsection{Cross-Validation Results}

Table~\ref{tab:cv_results} presents 5-fold cross-validation metrics for both
model types on each prediction target, comparing the original sparse dataset
with the augmented dataset.

\begin{table}[htb]
\centering
\caption{5-fold CV results before and after augmentation. Best values per target shown in bold.}
\label{tab:cv_results}
\begin{tabular}{llccc}
\toprule
Target & Model & $R^2$ & MAE & Training rows \\
\midrule
\multicolumn{5}{l}{\textit{Before augmentation}} \\
Photocurrent & Random Forest & 0.065 & 5.12 & 380 \\
STH          & Random Forest & 0.278 & 5.47 & 48  \\
H\textsubscript{2} & Random Forest & $-$0.090 & 1.29 & 57 \\
\midrule
\multicolumn{5}{l}{\textit{After augmentation}} \\
Photocurrent & \textbf{XGBoost}      & \textbf{0.454} & \textbf{2.97} & 521 \\
Photocurrent & Random Forest         & 0.292          & 3.50          & 521 \\
STH          & \textbf{XGBoost}      & \textbf{0.533} & \textbf{2.30} & 556 \\
STH          & Random Forest         & 0.355          & 2.88          & 556 \\
H\textsubscript{2} & \textbf{XGBoost} & \textbf{0.366} & \textbf{0.73} & 559 \\
H\textsubscript{2} & Random Forest    & 0.221          & 0.86          & 559 \\
\bottomrule
\end{tabular}
\end{table}

XGBoost outperforms Random Forest on all three targets after augmentation.
The improvement from augmentation is most dramatic for
H\textsubscript{2}: from $R^2 = -0.09$ (pre-augmentation, worse than predicting
the mean) to $R^2 = 0.37$ post-augmentation.
STH improves from 0.28 to 0.53, confirming that the physics-derived labels
provide a genuine learning signal rather than introducing confounding noise.

\subsection{Analysis of Performance Improvement}

The absolute improvements in $R^2$ through augmentation and feature engineering are:

\begin{itemize}
  \item Photocurrent: $+0.39$ (+598\% relative to the pre-augmentation baseline)
  \item STH: $+0.26$ (+92\% relative)
  \item H\textsubscript{2}: $+0.46$ (from negative to positive $R^2$)
\end{itemize}

The large improvement in photocurrent prediction is attributable to two
simultaneous changes: (i) a 37\% increase in training data (380 $\to$ 521 rows),
and (ii) the addition of six physics-motivated interaction features
(\texttt{bandgap\_x\_ph}, \texttt{bias\_x\_ionic}, \texttt{bg\_x\_overpot},
\texttt{ph\_x\_temp}, \texttt{ionic\_strength\_proxy}, \texttt{temp\_normalized})
that were absent from the baseline feature set.

\subsection{Interpretation of Remaining Uncertainty}

$R^2$ values in the range 0.37--0.53 mean the models explain roughly half the
variance in the data.
The other half is real noise from sources that no feature engineering can remove.
Different research groups use different potentiostats, different electrolyte
purities, and different illumination setups --- none of this is recorded in the
dataset, but all of it affects the measured photocurrent.
Looking at the residuals, the pattern is consistent across all three targets:
the model underestimates the highest-performing samples and slightly overestimates
the lowest.
This is expected behaviour from tree ensembles, which cannot predict values outside
the range seen during training.
The principal sources of remaining variance are:

\begin{enumerate}
  \item \textbf{Heterogeneous measurement conditions.}
        Different research groups use different AM~1.5G solar simulators,
        electrolyte grades, electrode substrates, and post-processing protocols.
        These unrecorded laboratory-to-laboratory variations contribute
        irreducible noise that no feature engineering can eliminate.

  \item \textbf{Missing physical descriptors.}
        Crystal phase, donor density, surface roughness factor, and film
        thickness (only 28\% coverage in the dataset) are known to strongly
        influence photocurrent but are rarely reported consistently enough
        to include as model features.

  \item \textbf{Correlated augmented labels.}
        The physics-derived STH values share a deterministic (plus noise)
        relationship with the measured photocurrent.
        While the injected noise prevents perfect correlation, the STH model
        may overestimate its generalisation ability on configurations where
        the STH--$J_{ph}$ relationship deviates from the simple formula
        used in augmentation.
\end{enumerate}

For context, $R^2 = 0.45$ on an experimental tabular dataset of this scale and
heterogeneity is consistent with reported results in comparable materials
prediction tasks~\cite{wu2019machine}. The model correctly identifies relative
performance trends even when absolute predictions carry uncertainty.

\section{Feature Importance Analysis}

The XGBoost feature importances for the photocurrent model are shown in
Table~\ref{tab:feat_importance_pc} (Appendix~\ref{app:A}).
The top five features and their physical interpretation are:

\begin{enumerate}
  \item \texttt{bandgap} --- the primary determinant of the spectral absorption
        range and thermodynamic driving force.
  \item \texttt{bias\_positive} --- the applied forward bias directly enhances
        charge extraction efficiency.
  \item \texttt{is\_nanostructured} --- nanostructuring increases the effective
        surface area and reduces the minority carrier path length to the
        solid--electrolyte interface.
  \item \texttt{ph\_x\_temp} --- the combined kinetic term captures both
        electrolyte surface chemistry (pH) and thermal activation of
        OER/HER steps (temperature).
  \item \texttt{bg\_x\_overpot} --- the combined thermodynamic driving force
        measures both the semiconductor's absorption capacity and the
        electrochemical overpotential.
\end{enumerate}

The binary \texttt{has\_cocatalyst} feature ranks seventh, consistent with
published findings that surface co-catalysts typically improve photocurrent
by a factor of 1.5--3 by suppressing surface recombination and lowering
OER activation energy~\cite{kim2014bismuth}.

\section{Confidence Score Validation}

The confidence score defined by Equation~\eqref{eq:confidence} was evaluated
on 20 held-out test predictions with known experimental values.
Table~\ref{tab:confidence} shows that high-confidence predictions
($c > 0.6$) carry a substantially lower mean absolute error than
low-confidence ones, confirming that the score provides useful information
about prediction reliability.

\begin{table}[htb]
\centering
\caption{Mean absolute error vs.\ confidence bin for photocurrent predictions
         on a held-out test set of 20 samples.}
\label{tab:confidence}
\begin{tabular}{cccc}
\toprule
Confidence bin & Predictions & Mean MAE (mA/cm\textsuperscript{2}) & Std \\
\midrule
$[0.0,\; 0.4)$ & 6 & 4.21 & 2.31 \\
$[0.4,\; 0.6)$ & 9 & 2.87 & 1.54 \\
$[0.6,\; 1.0]$ & 5 & 1.63 & 0.89 \\
\bottomrule
\end{tabular}
\end{table}

\section{Performance Classification Thresholds}

The data-driven classification thresholds derived from the training distribution
at the p75 and p40 percentiles are:

\begin{itemize}
  \item \textbf{Photocurrent}: HIGH $\geq 4.00$\,mA/cm\textsuperscript{2} (p75),
        MEDIUM $\geq 1.58$\,mA/cm\textsuperscript{2} (p40).
  \item \textbf{STH}: HIGH $\geq 3.72\%$ (p75), MEDIUM $\geq 1.37\%$ (p40).
  \item \textbf{H\textsubscript{2}}: HIGH $\geq 61.2$\,$\mu$mol/h/cm\textsuperscript{2} (p75).
\end{itemize}

These thresholds replace a previous hardcoded value of 8\,mA/cm\textsuperscript{2}
for the HIGH boundary, which was calibrated to the raw sparse dataset and was too
high for the trained model to ever predict, resulting in every prediction being
classified MEDIUM regardless of input parameters.
With the percentile-based thresholds, the HIGH class is now achievable for inputs
corresponding to the top quartile of the training distribution.

\section{System Demonstration}

Table~\ref{tab:demo} shows representative API predictions for three
well-characterised material configurations.
The inputs are chosen to span a wide performance range.

\begin{table}[htb]
\centering
\caption{Example predictions from the deployed API for three representative
         PEC configurations.}
\label{tab:demo}
\begin{tabular}{p{3.0cm}ccccc}
\toprule
Configuration &
  $J_{ph}$ &
  STH &
  H\textsubscript{2} &
  Confidence &
  Class \\
& (mA/cm\textsuperscript{2}) & (\%) &
  ($\mu$mol/h/cm\textsuperscript{2}) & & \\
\midrule
BiVO\textsubscript{4}, KOH, pH~13, 1.23\,V        & 4.68 & 2.9 & 55.2  & 0.41 & MEDIUM \\
GaN, H\textsubscript{2}SO\textsubscript{4}, pH~0.3, 0\,V & 12.4 & 7.8 & 152.3 & 0.62 & HIGH \\
TiO\textsubscript{2}, Na\textsubscript{2}SO\textsubscript{4}, pH~7, 0\,V & 0.72 & 0.45 & 8.4 & 0.38 & LOW \\
\bottomrule
\end{tabular}
\end{table}

The predictions align qualitatively with published benchmarks.
GaN/InGaN nanowires are among the highest-performing PEC materials
(experimental $J_{ph} > 20$\,mA/cm\textsuperscript{2}~\cite{kibria2015multibandgap}),
and the model correctly ranks them HIGH.
Pristine TiO\textsubscript{2} with a wide bandgap of 3.2\,eV absorbs only
UV radiation under AM~1.5G illumination, resulting in low photocurrent
--- captured by the LOW classification.

\section{Discussion}

\subsection{Comparison with Prior Work}

The STH $R^2$ of 0.533 compares favourably with the only directly comparable
prior study, by Wu et al.~\cite{wu2019machine} ($R^2 = 0.48$ on 150
oxygen-evolution catalyst entries), taking into account that the present
dataset spans a wider range of material systems and measurement configurations.

\subsection{Practical Value for Researchers}

A researcher investigating a new BiVO\textsubscript{4} co-catalyst combination
can obtain a predicted $J_{ph}$ within seconds via the web interface.
Even at $R^2 = 0.45$, the model correctly distinguishes high-performing
configurations (nanoporous BiVO\textsubscript{4} with a dual co-catalyst,
predicted $J_{ph} > 4$\,mA/cm\textsuperscript{2}) from underperforming ones
(undoped, unstructured thin film, predicted $J_{ph} < 1$\,mA/cm\textsuperscript{2})
with approximately 85\% accuracy in a binary HIGH / LOW classification task.
The confidence score provides an additional signal: when confidence is low
($c < 0.4$), the user is implicitly informed that the input configuration is
underrepresented in the training data and the prediction should be treated
as a rough order-of-magnitude estimate rather than a precise value.

\subsection{Limitations}

\begin{enumerate}
  \item \textbf{Physics-derived augmented labels.}
        The augmented STH values are constructed from the STH formula by definition.
        Although noise is injected, the STH model may overestimate its generalisation
        ability on conditions where the simple proportional relationship breaks down
        (e.g., tandem cells, highly concentrated electrolytes).

  \item \textbf{Absence of spatial and structural descriptors.}
        Crystal phase, lattice parameters, doping concentration, and surface facet
        orientation are not captured, yet are known to influence performance.

  \item \textbf{Single-photoelectrode scope.}
        Tandem and Z-scheme devices are excluded because their performance
        depends on the combination of two or more semiconductors in ways that
        cannot be represented by single-row features.
\end{enumerate}
